%%%%%%%%%%%%%%%%%%%%%%%%%%%%%%%%%%%%%%%%%%%%%%%%%%%%%%%%%%%
%% Front Matter — Chandrasekhar, "Hydrodynamic and
%%   Hydromagnetic Stability" (Clarendon Press, 1961)
%% Transcribed from scanned PDF pages 1–16
%%%%%%%%%%%%%%%%%%%%%%%%%%%%%%%%%%%%%%%%%%%%%%%%%%%%%%%%%%%

% ============================================================
% TITLE PAGE  (PDF p. 1)
% ============================================================
\begin{titlepage}
\centering
{\Huge\bfseries HYDRODYNAMIC\\AND\\HYDROMAGNETIC\\STABILITY\par}
\vspace{2cm}
{\large BY\par}
\vspace{0.5cm}
{\Large\bfseries S.~CHANDRASEKHAR\par}
\vspace{0.3cm}
{\normalsize UNIVERSITY OF CHICAGO\par}
\vfill
{\large OXFORD\par}
{\large AT THE CLARENDON PRESS\par}
{\large 1961\par}
\end{titlepage}

% ============================================================
% COPYRIGHT / PUBLISHER PAGE  (PDF p. 2)
% ============================================================
\thispagestyle{empty}
\noindent
{\itshape Oxford University Press, Amen House, London E.C.4}

\medskip
\noindent
\textsc{Glasgow \quad New York \quad Toronto \quad Melbourne \quad Wellington}\\
\textsc{Bombay \quad Calcutta \quad Madras \quad Karachi \quad Kuala Lumpur}\\
\textsc{Cape Town \quad Ibadan \quad Nairobi \quad Accra}

\bigskip
\begin{center}
\rule{3cm}{0.4pt}
\end{center}
\bigskip

\begin{center}
PRINTED IN GREAT BRITAIN
\end{center}

\cleardoublepage

% ============================================================
% PREFACE  (PDF p. 3, book p. v)
% ============================================================
\chapter*{PREFACE}
\addcontentsline{toc}{chapter}{Preface}
\markboth{PREFACE}{PREFACE}

\noindent
\textsc{In} this book I have tried to bring together into a coherent account
what I have learnt of hydrodynamic and hydromagnetic stability, with
my students and with my associates, during the past years.

I regret that even in a book of this size, I have not been able to do
justice to many aspects of the subject which could claim as much right
to attention as those that have received their due or ample share.
Perhaps the most serious omission is the absence of any reference to
viscous shear flow; but this is a large and a highly specialized field and
I have not felt myself competent to write about it. This argument is
indeed a general one: in the last analysis an author chooses to write
about only those matters in which he has some confidence of his understanding.

In the writing of this book I have received co-operation and assistance,
in generous measure, from so many that in order to be reasonably complete I have listed them separately.  I should, however, like to mention
here the extent of my obligation to Miss Donna~D. Elbert: in a real
sense this book is the outcome of our joint efforts over the years and
without her part there would have been no substance.

\hfill S.~C.

\cleardoublepage

% ============================================================
% ACKNOWLEDGEMENTS  (PDF pp. 4–5, book pp. vii–viii)
% ============================================================
\chapter*{ACKNOWLEDGEMENTS}
\addcontentsline{toc}{chapter}{Acknowledgements}
\markboth{ACKNOWLEDGEMENTS}{ACKNOWLEDGEMENTS}

\noindent
The following persons provided copies for many of the illustrations in this
book; and in several instances new computations or new experiments were
necessary.

\medskip
\begin{tabbing}
\hspace{1em}\= Dr.~F.~E.~Bisshopp (Figs.~5, 6, 9, and 57) \\
\> Dr.~R.~J.~Donnelly (Figs.~77, 78\,\textit{a}, \textit{b}, and 79\,\textit{a}, \textit{b}) \\
\> Drs.~R.~J.~Donnelly and D.~Fultz (Figs.~80, 82, 83, 84, 87, and 97) \\
\> Dr.~J.~R.~D.~Francis (Fig.~117) \\
\> Dr.~D.~Fultz (Figs.~20, 33, 35, and 64) \\
\> Mr.~I.~Gocol (Figs.~38\,\textit{a}, \textit{b}) \\
\> Dr.~B.~Lehnert (Fig.~46) \\
\> Dr.~Y.~Nakagawa (Figs.~30, 31, 43, 44, 45, 51, 52, 53, and 54) \\
\> Dr.~A.~H.~Nissan (Figs.~96\,\textit{a}, \textit{b}) \\
\> Dr.~W.~H.~Reid (Figs.~65, 66, 87, 88, 85, 99, 100, 108, and 113) \\
\> Drs.~W.~H.~Reid and D.~L.~Harris (Figs.~2, 3, 4, and 8) \\
\> Dr.~T.~L.~Silveston (Figs.~15\,\textit{a}, \textit{b}, 12, 14, and 17\,\textit{a}, \textit{b}) \\
\> Dr.~R.~K.~Trebus (Figs.~130 and 131) \\
\> Dr.~P.~O.~Vandervoort (Figs.~27, 28, 42, 71, 112, 114, 121, 135, and 136) \\
\> Dr.~G.~Veronis (Figs.~24\,\textit{a}, \textit{b}, 25\,\textit{a}, \textit{b}, and 26) \\
\> Mr.~J.~P.~Wright (Fig.~41)
\end{tabbing}

\medskip
\noindent
Permission was granted by the Editors of the following periodicals for the
reproduction of some of the illustrations:

\medskip
\begin{tabbing}
\hspace{1em}\= \textit{The Royal Society} (Figs.~11, 15, 16, 22, 34, 44, 50, 55, 68, 85, and 86) \\
\> \textit{The Journal of Fluid Mechanics} (Figs.~24\,\textit{a}, \textit{b}, 25\,\textit{a}, \textit{b}, and 26) \\
\> \textit{The Philosophical Magazine} (Fig.~117) \\
\> \textit{Proceedings of the Physical Society} (Figs.~133 and 134)
\end{tabbing}

\medskip
\noindent
Mr.~Norman R.~Wolfe inked all of the line drawings in the book.

\medskip
\noindent
The following persons read various parts of the completed manuscript and
made helpful suggestions.

\medskip
\begin{tabbing}
\hspace{1em}\= Dr.~W.~H.~Reid read Chapters~I--IX and contributed new information \\
\> \hspace{2em} included in Chapters~II, VII, VIII, IX, and~X \\
\> Dr.~P.~Ledoux read Chapters~I--VI \\
\> Dr.~M.~N.~Rosenbluth read Chapters~IX--XII
\end{tabbing}

\medskip
\noindent
The experimental investigations of Drs.~D.~Fultz, Y.~Nakagawa, and R.~J.
Donnelly in close association with the theoretical developments have been
a source of constant encouragement; much of the experimental work summarized in this book is theirs.

% --- Page viii ---

Dr.~Norman R.~Lebovitz read the entire book in proofs and helped to
eliminate many errors.

Miss Donna~D.~Elbert carried out the relevant numerical calculations for
most of the problems treated in this book; she is responsible for the numerical
information included in all of the tables with the exception of Tables~I--VI,
X, XXIII--XXX, XXXVI--XXXIX, XLVII, XLIX, LXIV--LXVII, and
LXX.

And I am grateful to all of these persons.

\bigskip

\noindent
The theoretical investigations included in this book have in large part been
supported by contracts with the Geophysics Research Directorate of the Air
Force Cambridge Research Center, Air Research and Development Command
(Contracts AF~19(604)--299 and AF~19(604)--2046).

And the experimental investigations carried out at the Enrico Fermi
Institute for Nuclear Studies (University of Chicago) have been supported
by contracts with the Office of Naval Research (Contracts N6ori-02056 and
Nonr-2121(20)).

\bigskip

\noindent
And finally, I am grateful to the Clarendon Press for bringing to this book
that excellence of craftsmanship and typography which is characteristic of
all of their work.

\hfill S.~C.

\cleardoublepage

% ============================================================
% TABLE OF CONTENTS — transcribed as comments for reference
% (LaTeX will auto-generate the actual TOC via \tableofcontents)
% ============================================================
%
% The original TOC spans PDF pages 6–16 (book pages ix–xix).
% It is transcribed below as comments so that the structure can
% be verified against the auto-generated LaTeX TOC.
%
% ------------------------------------------------------------------
% CONTENTS
% ------------------------------------------------------------------
%
% I. BASIC CONCEPTS                                                  1
%    1. Introduction                                                 1
%    2. Basic concepts                                               1
%    3. The analysis in terms of normal modes                        3
%    4. Non-dimensional numbers                                      6
%    BIBLIOGRAPHICAL NOTES                                           7
%
% II. THE THERMAL INSTABILITY OF A LAYER OF
%     FLUID HEATED FROM BELOW                                        9
%     1. THE B\'ENARD PROBLEM
%    5. Introduction                                                 9
%    6. The nature of the physical problem                           9
%    7. The basic hydrodynamic equations                            10
%       (a) The equation of continuity                              10
%       (b) The equations of motion                                 11
%       (c) The rate of viscous dissipation                         14
%       (d) The equation of heat conduction                         14
%    8. The Boussinesq approximation                                16
%    9. The perturbation equations                                  18
%       (a) The boundary conditions                                 21
%   10. The analysis into normal modes                              21
%       (a) The solution for the horizontal components of the velocity  24
%   11. The principle of the exchange of stabilities                24
%   12. The equations governing the marginal state and the reduction to
%       a characteristic value problem                              25
%   13. The variational principles                                  27
%       (a) The first variational principle                         27
%       (b) The second variational principle                        31
%   14. The thermodynamic significance of the variational principle  34
%   15. Exact solutions of the characteristic value problem         34
%       (a) The solution for two free boundaries                    35
%       (b) The solution for two rigid boundaries                   36
%           (i) The even solutions                                  37
%           (ii) The odd solutions                                  39
%       (c) The solution for one rigid and one free boundary        42
%       (d) Summary of the results for the three cases              43
%   16. The cell patterns                                           43
%       (a) Rolls                                                   44
%       (b) Rectangular and square cells                            44
%       (c) Hexagonal cells                                         47
%       (d) Triangular cells                                        50
%       (e) More general cell patterns                              51
%   17. The variational solution                                    53
%   18. Experiments on the onset of thermal instability in fluids   59
%       (a) B\'enard's experiments                                    60
%       (b) The Schmidt--Milverton principle for detecting the onset of
%           thermal instability                                     61
%       (c) The precision experiments of Silveston                  63
%       (d) Observations by optical methods                         65
%   BIBLIOGRAPHICAL NOTES                                           68
%
% III. THE THERMAL INSTABILITY OF A LAYER OF
%      FLUID HEATED FROM BELOW                                      76
%      2. THE EFFECT OF ROTATION
%   19. Introduction                                                76
%   20. The theorems of Helmholtz and Kelvin                        77
%   21. The equations of hydrodynamics in a rotating frame of reference  80
%   22. The Taylor--Proudman theorem                                 82
%   23. The propagation of waves in a rotating fluid                83
%   24. The problem of thermal instability in a rotating fluid: general
%       considerations                                              87
%   25. The perturbation equations                                  87
%       (a) The analysis into normal modes                          90
%   26. The case when instability sets in as stationary convection. A vari-
%       ational principle                                           91
%   27. Solutions for the case when instability sets in as stationary con-
%       vection                                                     94
%       (a) The solution for two free boundaries                    94
%       (b) The solution for two rigid boundaries                   97
%       (c) The solution for one rigid and one free boundary       101
%       (d) The origin of the $\Pi$-set                            104
%   28. The motions in the horizontal planes and the cell patterns at the
%       onset of instability as stationary convection              106
%       (a) Rolls, rectangles, and squares                         106
%       (b) Hexagons                                               111
%   29. The limiting behaviour of the streamlines for $\mathscr{T} \to \infty$  113
%   29. On the onset of convection as overstability. The solution for the
%       case of two free boundaries                                114
%       (a) The nature of the roots of the characteristic equation (215)  121
%   30. On a method of discriminating the character of the marginal state.
%       A general variational principle                            123
%       (a) The variational principle                              124
%
% --- Page x (PDF p. 7) ---
%
%   31. The onset of convection as overstability: the solution for other
%       boundary conditions                                        126
%       The case $p = 0$                                           126
%   32. Thermodynamic significance of the variational principles   128
%       (a) The case when the marginal state is stationary         129
%       (b) The case when the marginal state is oscillatory        132
%       (c) The case when $\Omega$ and $g$ act in different directions  133
%   33. Experiments on the onset of thermal instability in rotating fluids 134
%       (a) Experiments with water                                 134
%       (b) Experiments with mercury                               136
%   BIBLIOGRAPHICAL NOTES                                          138
%
% IV. THE THERMAL INSTABILITY OF A LAYER OF
%     FLUID HEATED FROM BELOW                                      140
%     3. THE EFFECT OF A MAGNETIC FIELD
%   34. Hydromagnetics                                             140
%   35. The basic equations of hydromagnetics                      140
%   36. The equation of motion governing the magnetic field and some of
%       its consequences                                           143
%       (a) The decay of a magnetic field in the absence of fluid motions.
%           The Joule dissipation                                  145
%       (b) The case when there are motions and the conductivity is infinite 147
%           (i) Convection theorems                                148
%       (c) The transformation of magnetic energy into kinetic energy
%           and conversely                                         149
%       (d) The general energy equation                            150
%   39. The Alfv\'en waves                                          150
%       (a) The case when $\epsilon = \eta = 0$                    153
%       (b) The effects of finite viscosity and resistivity        154
%   40. Some special solutions of the hydrodynamic equations. The
%       analogues of the Taylor--Proudman theorem                  157
%       (a) The equipartition solutions                            158
%       (b) Force-free fields                                      158
%       (c) The analogue of the Taylor--Proudman theorem            158
%   41. The problem of thermal instability in the presence of a magnetic
%       field: general considerations                              159
%   42. The perturbation equations                                 160
%       (a) The boundary conditions                                162
%       (b) The analysis into normal modes                         163
%       (c) The solution for the horizontal components of the velocity and
%           the magnetic field                                     164
%   43. The case when instability sets in as stationary convection. A vari-
%       ational principle                                          166
%       (a) A variational principle                                166
%       (b) The thermodynamic significance of the variational principle  167
%
% --- Page xi (PDF p. 8) ---
%
%   44. Solutions for the case when instability sets in as stationary con-
%       vection                                                    170
%       (a) The solution for two free boundaries                   170
%       (b) The solution for two rigid boundaries                  172
%       (c) The solution for one rigid and one free boundary       174
%       (d) Cell patterns                                          176
%   45. The origin of the $\epsilon^{*}Q$ law and an invariant     178
%       (a) An invariant                                           179
%   46. On the onset of convection as overstability                181
%   47. The case when $H$ and $g$ act in different directions      186
%   48. Experiments on the inhibition of thermal convection by a magnetic
%       field                                                      189
%   BIBLIOGRAPHICAL NOTES                                          193
%
% V. THE THERMAL INSTABILITY OF A LAYER OF
%    FLUID HEATED FROM BELOW                                       196
%    4. THE EFFECT OF ROTATION AND MAGNETIC FIELD
%   49. The like and the contrary effects of rotation and magnetic field on
%       fluid behaviour                                            196
%   50. The propagation of hydromagnetic waves in a rotating fluid 199
%   51. The perturbation equations                                 202
%   52. The case when instability sets in as stationary convection 202
%       (a) The solution for a case of two free boundaries         203
%   53. The case when instability sets in as overstability         204
%       (a) An approximate solution applicable to liquid metals    205
%       (b) Numerical results for mercury                          209
%   54. Experiments on the onset of thermal instability in the presence of
%       rotation and magnetic field                                211
%       (a) The results on the critical Rayleigh number and on the manner
%           of the onset of instability                            212
%       (b) Optical observations and the discontinuous variation of the cell
%           dimension with the strength of the magnetic field      217
%   BIBLIOGRAPHICAL NOTES                                          219
%
% VI. THE ONSET OF THERMAL INSTABILITY IN FLUID
%     SPHERES AND SPHERICAL SHELLS                                 220
%   55. Introduction                                               220
%   56. The perturbation equations                                 222
%       (a) The operator $\mathscr{L}$                             222
%       (b) The analysis into normal modes                         223
%       (c) Boundary conditions                                    224
%       (d) The velocity field                                     225
%   57. The validity of the principle of the exchange of stability for the
%       case $\beta = \text{constant}$ and $\gamma = \text{constant}$  226
%   58. A variational principle for the case when $\beta$ and $\gamma$ are constants 229
%       (a) The thermodynamic significance of the variational principle  230
%
% --- Page xii (PDF p. 9) ---
%
%   59. On the onset of thermal instability in a fluid sphere      231
%       (a) The cell patterns                                      234
%   60. On the onset of thermal instability in spherical shells    237
%       (i) Free surfaces at $r = 1$ and $r = q$                  237
%       (ii) A rigid surface at $r = r_0$ and a free surface at $r = 1$  241
%       (iii) A free surface at $r = r_0$ and a rigid surface at $r = 1$  242
%       (iv) Rigid surfaces at $r = 1$ and $r = q$                242
%       (a) The case $h = a = 1$                                   245
%       (b) The case $h(r) = 1$                                    247
%       (c) The case $\alpha(r) = 1$                               248
%   61. On the effect of rotation on the onset of thermal instability in a fluid
%       sphere. The formulation of the problem                     251
%       (a) The representation of an axisymmetric solenoidal vector field 252
%       (b) The perturbation equations                             254
%       (c) The boundary conditions                                256
%       (d) The variational principle                              257
%       (e) The thermodynamic significance of the variational principle  259
%   62. The effect of rotation on the onset of stationary convection in a
%       fluid sphere                                               260
%   63. Some remarks on geophysical applications                   266
%   BIBLIOGRAPHICAL NOTES                                          268
%
% VII. THE STABILITY OF COUETTE FLOW                               271
%   64. Introduction                                               271
%   65. The physical problem                                       272
%   66. Rayleigh's criterion                                       273
%   67. Analytical discussion of the stability of inviscid Couette flow  277
%       (a) The equations in terms of the Lagrangian displacement  278
%       (b) The case $m = 0$                                       280
%       (c) The case $m \neq 0$                                    281
%   68. The periods of oscillation of a rotating column of liquid  282
%       (a) The case $\Omega = \text{constant}$                    285
%           (i) The case $\eta = 0$                                285
%           (ii) The case $m = 0$                                  286
%       (b) The case $\Omega = A + B/r^2$ and $m = 0$             286
%           (i) The solution for a narrow gap                      289
%           (ii) The formal solution for a wide gap                292
%   69. On viscous Couette flow                                    292
%   70. The perturbation equations                                 295
%       (a) The stability of the flow for $\mu > \eta^2$          296
%   71. The solution for the case of a narrow gap when the marginal state is
%       stationary                                                 298
%       (a) The solution of the characteristic value problem for the case
%           $d = 0$                                                300
%
% --- Page xiii (PDF p. 10) ---
%
%       (b) Numerical results                                      303
%       (c) An alternative method of solution                      307
%       (d) An approximate solution for $\mu = 1$                 309
%       (e) The asymptotic behaviour for $|\eta| \to 0$           313
%   72. On the principle of the exchange of stabilities            314
%   73. The solution for a wide gap when the marginal state is stationary 318
%       (a) The characteristic equation                            319
%       (b) Numerical results for the case $\nu = \frac{1}{2}$   320
%   74. Experiments on the stability of viscous flow between rotating
%       cylinders                                                  330
%       (a) The determination of the critical Taylor numbers, for the case
%           $\mu = 0$, by torque measurements                     331
%       (b) Results of the experiments with the narrow gap         332
%       (c) Results of the experiments with the wide gap ($\nu = \frac{1}{2}$) 332
%       (d) The dependence of the critical Taylor number on $\Omega_1/\Omega_2$. The
%           results of visual and photographic observations        333
%       (e) Observations on the wave numbers of the disturbance
%           maintained at marginal stability                       335
%       (f) Comparison between the measured and the predicted
%           $(T_a/T_c)$-relations                                 337
%   BIBLIOGRAPHICAL NOTES                                          339
%
% VIII. THE STABILITY OF MORE GENERAL FLOWS
%       BETWEEN COAXIAL CYLINDERS                                  343
%   75. Introduction                                               343
%   76. The stability of viscous flow in a curved channel          344
%       (a) The perturbation equations                             344
%       (b) The solution of the characteristic value problem for the case
%           $\sigma = 0$                                           348
%       (c) Numerical results                                      350
%   77. The stability of viscous flow between rotating cylinders when a
%       transverse pressure gradient is present                    350
%       (a) The perturbation equations for the case $(R_2 - R_1)(R_2'/R_1 - R_1')$ 350
%       (b) The solution of the characteristic value problem for the case
%           $\sigma = 0$ and $\mu = 0$                             352
%       (c) The physical interpretation of the results             353
%       (d) Comparison with experimental results                   355
%   78. The stability of inviscid flow between coaxial cylinders when an axial
%       pressure gradient is present                               361
%       (a) The case of a pure axial flow                          361
%       (b) The general case when rotation is also present         366
%           A variational principle for $\epsilon$                 368
%           (i) A variational principle for $N^2$                  369
%           (ii) The criterion for stability                       369
%
% --- Page xiv (PDF p. 11) ---
%
%   79. The stability of viscous flow between rotating coaxial cylinders
%       when an axial pressure gradient is present                 371
%       (a) The perturbation equations                             372
%       (b) The reduction to the case of a narrow gap              373
%       (c) An approximate solution of the characteristic value problem
%           for the case $\mu > 0$                                 375
%       (d) Comparison with experimental results                   377
%   BIBLIOGRAPHICAL NOTES                                          379
%
% IX. THE STABILITY OF COUETTE FLOW IN
%     HYDROMAGNETICS                                               381
%   80. The equations of hydromagnetics in cylindrical polar coordinates  382
%   81. The stability of non-dissipative Couette flow when a magnetic field
%       parallel to the axis is present                            383
%       (a) The case $m = 0$                                       386
%       (b) The case $m \neq 0$                                    387
%   82. The periods of oscillation of a rotating column of liquid when a
%       magnetic field is impressed in the direction of the axis   388
%       (a) The case $\Omega$ is constant                          388
%       (b) The case $m = 0$ and $\Omega = A + B/r^2$, the stabilizing effect of a
%           magnetic field in case of narrow gaps                  390
%   83. The stability of non-dissipative Couette flow when a current flows
%       parallel to the axis                                       391
%       (a) The case $m = \infty$                                  393
%       (b) The case $m = 0$                                       394
%   84. The stability of non-dissipative Couette flow when an axial and a
%       transverse magnetic field are present                      396
%       (a) The case $\Omega = 0$                                  397
%   85. The stability of dissipative Couette flow in hydromagnetics. The
%       perturbation equations                                     399
%       (a) The boundary conditions                                400
%       (b) The equations governing the marginal state when the onset of
%           instability is as a stationary secondary flow          401
%       (c) The reduction to the case of a narrow gap              404
%   86. The solution of the characteristic value problem for the case
%       $p = 0$                                                    406
%       (a) A variational principle                                406
%       (b) The solution of the characteristic value problem. The secular
%           determinants                                           407
%           The case of non-conducting walls                       408
%           The case of conducting walls                           410
%       (c) Numerical results                                      411
%       (d) The asymptotic behaviour for $Q \to \infty$            413
%   87. The solution of the characteristic value problem in the general case 415
%       (a) The case $\mu = -1$                                    417
%
% --- Page xv (PDF p. 12) ---
%
%   88. The stability of dissipative flow in a curved channel in the presence
%       of an axial magnetic field                                 422
%       (a) The equations governing the marginal state when the onset of
%           instability is as a stationary secondary flow          423
%       (b) The stability of a pure pressure maintained flow       425
%   89. Experiments on the stability of viscous flow between rotating
%       cylinders when an axial magnetic field is present          426
%   BIBLIOGRAPHICAL NOTES                                          428
%
% X. THE STABILITY OF SUPERPOSED FLUIDS: THE
%    RAYLEIGH--TAYLOR INSTABILITY                                  428
%   90. Introduction                                               428
%   91. The character of the equilibrium of a stratified heterogeneous
%       fluid. The perturbation equations                          430
%       (a) Allowance for surface tension at interfaces between fluids  430
%   92. The inviscid case                                          435
%       (a) The case of two uniform fluids of constant density separated
%           by a horizontal boundary                               435
%       (b) The case of exponentially varying density              438
%   93. A general variational principle                            441
%       The variational principle                                  442
%   94. The case of two uniform viscous fluids separated by a horizontal
%       boundary                                                   443
%       (a) The case $\rho_1 = \rho_2$                             444
%       (b) The modes of maximum instability for the case $\rho_1 = \rho_2$,
%           $\rho_1 > \rho_2$ and $\sigma = 0$                    446
%       (c) The effect of surface tension on the unstable modes for $\alpha_1 = \alpha_2$,
%           $\rho_1 > \rho_2$                                     448
%       (d) The manner of decay in the case $\rho_1 = \rho_2$, $\rho_1 < \rho_2$, and $\mathscr{S} = 0$  448
%       (e) Gravity waves                                          450
%   95. The effect of rotation                                     453
%   96. The case of two uniform fluids separated by a horizontal boundary 455
%       The case of exponentially varying density                  456
%   97. The effect of a vertical magnetic field                    457
%       (a) Two uniform fluids separated by a horizontal boundary: the
%           unstable case                                          459
%       (b) Two uniform fluids separated by a horizontal boundary: the
%           stable case                                            462
%   97. The effect of a horizontal magnetic field                  464
%   98. The oscillations of a viscous liquid globe                 467
%       (a) The perturbation equations and their solution          467
%       (b) The solution for the inviscid case: the Kelvin modes   469
%       (c) The solution for the general case                      470
%       (d) The boundary conditions and the characteristic equation  470
%       (e) The manner of decay of the Kelvin modes. Numerical results 472
%
% --- Page xvi (PDF p. 13) ---
%
%   99. The oscillations of a viscous liquid drop                  475
%       (a) The solution for the inviscid case                     476
%       (b) The solution for the general case                      476
%   BIBLIOGRAPHICAL NOTES                                          477
%
% XI. THE STABILITY OF SUPERPOSED FLUIDS: THE
%     KELVIN--HELMHOLTZ INSTABILITY                                481
%  100. The perturbation equations                                 481
%  101. The case of two uniform fluids in relative horizontal motion
%       separated by a horizontal boundary                         483
%       (a) The case when surface tension is absent                484
%  102. The stabilizing effect of surface tension                  485
%  102. The effect of a continuous variation of $U$ on the development of the
%       Kelvin--Helmholtz instability                              487
%  103. The Kelvin--Helmholtz instability in a fluid in which both $\rho$ and $U$
%       are continuously variable                                  488
%       (a) Analytical results for the case $\rho = \rho_0 e^{-\alpha z}$ and $U = U_0 + b z$ 491
%  104. An example of the instability of a shear layer in an unbounded
%       heterogeneous inviscid fluid                               495
%  105. The effect of rotation                                     499
%       (a) The case of two uniform fluids in relative horizontal
%           motion separated by a horizontal boundary              499
%  105. The discussion of the characteristic equation              501
%  106. The effect of a horizontal magnetic field                  507
%       (a) The effect of a magnetic field in the direction of streaming  508
%       (b) The case of two uniform fluids in relative horizontal
%           motion separated by a horizontal boundary              510
%       (c) The effect of a magnetic field transverse to the direction of
%           streaming                                              513
%   BIBLIOGRAPHICAL NOTES                                          515
%
% XII. THE STABILITY OF JETS AND CYLINDERS                         515
%  107. Introduction                                               515
%  108. The gravitational instability of an infinite cylinder      516
%       (a) The origin of the gravitational instability and an alternative
%           method of determining the characteristic frequencies   520
%  109. The effect of viscosity on the gravitational instability of an
%       infinite cylinder                                          525
%       (a) The case when the effects of viscosity are dominant    527
%       (b) The general case                                       530
%  110. The effect of a uniform axial magnetic field on the gravitational
%       instability of an infinite cylinder                        531
%       (a) The nature of the stabilizing effect of the axial magnetic field 534
%  111. The capillary instability of a liquid jet                  537
%       (a) The origin of the capillary instability                539
%
% --- Page xvii (PDF p. 14) ---
%
%       (b) The capillary instability of a hollow jet              539
%       (c) The effect of viscosity                                540
%  112. The effect of a uniform axial magnetic field on the capillary in-
%       stability of a liquid jet                                  543
%       (a) The effect of finite electrical conductivity           549
%       (b) The case of high resistivity                           549
%       (c) The general case                                       551
%  113. The stability of the simplest solution of the equations of hydro-
%       magnetics                                                  553
%       (a) An example                                             554
%  114. The effect of fluid motions on the stability of helical magnetic fields 556
%       (a) The case $f = 0$                                       558
%       (b) The case $f = 1$                                       561
%       (c) The general case                                       563
%  115. The stability of the pinch                                 566
%       (a) On stable pinch configurations                         571
%   BIBLIOGRAPHICAL NOTES                                          573
%
% XIII. GRAVITATIONAL EQUILIBRIUM AND
%       GRAVITATIONAL INSTABILITY                                  577
%  116. Introduction                                               577
%  117. The virial theorem                                         577
%       (a) Some definitions and relations                         577
%       (b) The general form of the virial theorem                 579
%       (c) The virial theorem for equilibrium configurations      581
%  118. The virial theorem for small oscillations about equilibrium 583
%       (a) A characteristic equation for determining the periods of oscil-
%           lation                                                 585
%  119. The gravitational instability of an infinite homogeneous medium;
%       Jeans's criterion                                          589
%  120. The effect of a uniform rotation and a uniform magnetic field on
%       Jeans's criterion                                          592
%       (a) The effect of a uniform rotation                       592
%       (b) The effect of a uniform magnetic field                 594
%       (c) The effect of the simultaneous presence of rotation and magnetic
%           field                                                  596
%   BIBLIOGRAPHICAL NOTES                                          598
%
% XIV. A GENERAL VARIATIONAL PRINCIPLE                             600
%  121. Introduction                                               600
%  122. A variational principle for treating the stability of hydromagnetic
%       systems                                                    601
%       (a) The effect of viscosity                                605
%  123. Extension of the variational principle to allow for compressibility 608
%   BIBLIOGRAPHICAL NOTES                                          608
%
% --- Page xviii (PDF p. 15) ---
%
% APPENDIX I. INTEGRAL RELATIONS GOVERNING STEADY
%    CONVECTION                                                    609
%  124. Introduction                                               609
%  125. The integral relations                                     609
%       (a) The integral relations                                 610
%  126. The amplitude of the steady convection past marginal stability 613
%   BIBLIOGRAPHICAL NOTES                                          616
%
% APPENDIX II. THE VARIATIONAL FORMULATION OF THE
%    PROBLEM CONSIDERED IN CHAPTER V                               617
%  127. The general equations governing the energy balance         617
%   BIBLIOGRAPHICAL NOTE                                           621
%
% APPENDIX III. TOROIDAL AND POLOIDAL VECTOR FIELDS                622
%  128. A general characterization of solenoidal vector fields. The funda-
%       mental basis                                               622
%  129. The orthogonality properties of the basic toroidal and poloidal
%       fields                                                     625
%   BIBLIOGRAPHICAL NOTES                                          626
%
% APPENDIX IV. VARIATIONAL METHODS BASED ON ADJOINT
%    DIFFERENTIAL SYSTEMS                                          627
%  130. Adjunctions of differential systems. An example            627
%  131. The dual relationship and the variational principle        629
%   BIBLIOGRAPHICAL NOTES                                          633
%
% APPENDIX V. ORTHOGONAL FUNCTIONS WHICH SATISFY FOUR
%    BOUNDARY CONDITIONS                                           634
%  132. Introduction                                               634
%  133. Functions suitable for problems with plane boundaries      634
%       (a) Functions involving $C_n(\zeta)$ and $S_n(\zeta)$     636
%  134. Functions suitable for problems with cylindrical and spherical
%       boundaries                                                 638
%       (a) The normalization integral                             638
%       (b) The cylinder functions of order 1                      640
%       (c) The spherical functions of half-odd integral orders    640
%   BIBLIOGRAPHICAL NOTES                                          642
%
% SUBJECT INDEX                                                    643
%
% INDEX OF DEFINITIONS                                             653
%
% --- Page xix (PDF p. 16) ---
%
% The following figures are printed as half-tone plates and will be found at the
% places given below:
%
% Fig. 1 facing p. 10; Figs. 15, 16, 17 between pp. 70-71; Fig. 20 facing p. 65;
% Figs. 30, 31 facing p. 136; Fig. 44 facing p. 190; Fig. 46 facing p. 193; Fig. 55
% facing p. 218; Fig. 64 facing p. 267; Fig. 68 facing p. 294; Figs. 82, 83, 84 between
% pp. 336-337; Fig. 96 facing p. 360; Fig. 117 facing p. 448.
%
% ------------------------------------------------------------------

\tableofcontents
